\SourceLeaf{31}kleinen Zeittheilchen $\mathrm{d}t$, in Folge der während dieser Zeit auf sie wirkenden Kräfte und der zur Zeit $t$ erlangten Geschwindigkeiten und Richtungen, einnehmen würden, falls sie alle vollkommen frei wären. Die wirklichen Plätze $c,c',c''$ u. s. w. werden dann diejenigen sein, für welche, unter allen mit den Bedingungen des Systems vereinbaren, $m(bc)^2+m'(b'c')^2+m''(b''c'')^2$ u. s. w. ein Minimum wird.

Das Gleichgewicht ist offenbar nur ein einzelner Fall des allgemeinen Gesetzes, und die Bedingung dafür, dass
\[
m(ab)^2+m'(a'b')^2+m''(a''b'')^2\ \text{u. s. w.}
\]
selbst ein Minimum sei, oder dass das Beharren des Systems im Zustande der Ruhe, der freien Bewegung der einzelnen Punkte näher liege, als jedes mögliche Heraustreten aus demselben.

Die Ableitung unsers Princips aus dem oben angeführten geschieht leicht auf folgende Art.

Die auf den materiellen Punkt $m$ wirkende Kraft ist offenbar zusammengesetzt, erstens aus einer, die, in Verbindung mit der zur Zeit $t$ Statt habenden Geschwindigkeit und Richtung, ihn in der Zeit $\mathrm{d}t$ von $a$ nach $c$ führt, und in einer zweiten, die ihn in derselben Zeit aus der Ruhe in $c$, durch $cb$ führen würde, wenn man den Punkt als frei betrachtet. Dasselbe gilt von den andern Punkten. Nach \textsc{Dalembert}’s Princip müssen demnach die Punkte $m,m',m''$ u. s. w., unter alleiniger Wirkung der zweiten Kräfte, nach $cb,c'b',c''b''$ u. s. w., in den Plätzen $c,c',c''$ u. s. w. vermöge der Bedingungen des Systems, im Gleichgewicht sein.

\expandafter\def\csname HistoricalNoteText3\endcsname{Der gewöhnliche Ausdruck setzt stillschweigend solche Bedingungen voraus, dass die jeder möglichen Bewegung entgegengesetzte gleichfalls möglich sei, wie z. B. dass ein Punkt auf einer bestimmten Fläche zu bleiben genöthigt, dass die Entfernung zweier Punkte von einander unveränderlich sei u. dgl. Allein dies ist eine unnöthige und der Natur nicht immer angemessene Beschränkung. Die Oberfläche eines undurchdringlichen Körpers zwingt einen auf ihr befindlichen materiellen Punkt nicht, auf ihr zu bleiben, sondern verwehrt ihr bloss das Austreten auf die Eine Seite; ein gespannter, nicht ausdehnbarer aber biegsamer Faden zwischen zwei Punkten macht nur die Zunahme, nicht die Abnahme der Entfernung unmöglich u. s. w. Warum wollten wir also das Gesetz der virtuellen Geschwindigkeiten nicht lieber gleich anfangs so ausdrücken, dass es \textit{alle} Fälle umfasst?}
Nach dem Princip der virtuellen Geschwindigkeiten erfordert dies Gleichgewicht, dass die Summe der Producte aus je drei Factoren, nemlich jeder der Massen $m,m',m''$ u. s. w., den Linien $cb,c'b'\ c''b''$ u. s. w., und irgend welchen auf letztere resp. projicirten, vermöge der Bedingungen des Systems möglichen Bewegungen jener Punkte, immer $=0$ sei, wie man es gewöhnlich ausspricht\HistoricalNote{3}
